\documentclass[11pt]{article}

\usepackage{amsmath,amssymb,amsfonts}
\usepackage[margin=1in]{geometry}
\usepackage{hyperref}
\usepackage{bm}
\usepackage{booktabs}

\hypersetup{colorlinks=true, linkcolor=blue, citecolor=blue, urlcolor=blue}

\newcommand{\vw}{\bm{v}}
\newcommand{\ee}{\bm{e}}
\newcommand{\KL}{D_{\mathrm{KL}}}
\newcommand{\PMI}{\mathrm{PMI}}
\newcommand{\BoW}{\mathrm{BoW}}

\title{Nearest-Neighbor Perturbation of the Gaussian\\Bag-of-Words Embedding Theory}
\author{}
\date{}

\begin{document}
\maketitle

\begin{abstract}
We consider a perturbation of the permutation-invariant bag-of-words (BoW) energy model for word co-occurrence, in which a nearest-neighbor interaction $\lambda\, \ee_{y_p} \cdot \ee_{y_{p+1}}$ between adjacent token embeddings is added to the Hamiltonian. The perturbation is self-referential: it is built from the same embedding dot products that the unperturbed Gaussian theory determines. We formulate the resulting self-consistency problem using the Cornwall--Jackiw--Tomboulis (CJT) 2PI effective action, derive the perturbative corrections to the embedding geometry, and identify the gap equation that determines the dressed propagator. Several open questions are posed regarding the sign of the coupling, the expansion base, and the role of the connected 4-point function.
\end{abstract}

\tableofcontents

\section{The unperturbed theory}\label{sec:unperturbed}

The Gaussian bag-of-words model assigns to a window of $L$ tokens $(y_1, \ldots, y_L)$ the probability
\begin{equation}\label{eq:H0}
    p_{\BoW}(y_1, \ldots, y_L) \propto \exp\!\Big(-\sum_{i,j} G_{ij}\, n_i\, n_j\Big) \equiv e^{-H_0}\,,
\end{equation}
where $n_i = \sum_{p=1}^L \delta_{i, y_p} - P_i$ is the count fluctuation of token $i$ and $G_{ij}$ is the coupling matrix. This distribution is invariant under $S_L$ (permutations of positions). The central result of the unperturbed theory is that the optimal embedding vectors $\ee_i^{(0)}$ satisfy
\begin{equation}\label{eq:bow-cov}
    \ee_i^{(0)} \cdot \ee_j^{(0)} = \frac{\Sigma_{ij}}{L\, P_i\, P_j}\,,
\end{equation}
where $\Sigma_{ij} = \langle n_i\, n_j \rangle - \langle n_i \rangle \langle n_j \rangle$ is the covariance of the count fluctuations. The dot product encodes the \emph{normalized covariance}---a linear measure of association between tokens. This is distinct from the pointwise mutual information $\PMI(i,j) = \log[p(i,j)/(p(i)p(j))]$, which is a logarithmic measure; the two are related but not equal. (The Arora et al.\ RAND-WALK model yields $\ee_i \cdot \ee_j \propto \PMI(i,j)$ under different assumptions; the present variational derivation gives the covariance directly.)

\section{The nearest-neighbor perturbation}\label{sec:perturbation}

\subsection{Definition}

We add a nearest-neighbor interaction in position space:
\begin{equation}\label{eq:V}
    V = \lambda \sum_{p=1}^{L-1} \ee_{y_p} \cdot \ee_{y_{p+1}}\,,
\end{equation}
where $\ee_{y_p}$ is the embedding of the token at position $p$ and $\lambda$ is the coupling strength. The full Hamiltonian is
\begin{equation}\label{eq:full-H}
    \boxed{\;H = \underbrace{\sum_{i,j} G_{ij}\, n_i\, n_j}_{H_0\;\text{(BoW)}} + \underbrace{\lambda \sum_{p=1}^{L-1} \ee_{y_p} \cdot \ee_{y_{p+1}}}_{V\;\text{(nearest-neighbor)}}\;}
\end{equation}

At $\lambda = 0$: pure bag-of-words (permutation-invariant). At $\lambda \neq 0$: the sum $\sum_p \ee_{y_p} \cdot \ee_{y_{p+1}}$ depends on the ordering of tokens, breaking the $S_L$ symmetry.

\subsection{Self-referential structure}

The perturbation $V$ is built from the embedding dot products $\ee_a \cdot \ee_b$, which are the same objects that the theory is meant to determine. This is analogous to a $\phi^4$ field theory in which the vertex coupling is constructed from the propagator---a self-referential or ``bootstrap'' structure. The appropriate formalism is the 2PI effective action (Section~\ref{sec:2PI}).

\subsection{Sign of $\lambda$}

The physical interpretation depends on the sign:
\begin{itemize}
    \item $\lambda < 0$: the nearest-neighbor term \emph{lowers} the energy when adjacent embeddings are aligned. This favors sequences where semantically similar tokens appear next to each other---reinforcing co-occurrence structure.
    \item $\lambda > 0$: the term \emph{raises} the energy for aligned adjacent tokens. This favors sequences where adjacent tokens are dissimilar---capturing the empirical fact that adjacent words tend to be syntactically complementary (e.g., determiner--noun, verb--object) rather than semantically identical.
\end{itemize}
The choice of sign is a modeling decision that determines whether the perturbation captures \emph{semantic} attraction or \emph{syntactic} complementarity between neighbors.

\section{Perturbative expansion}\label{sec:perturbative}

\subsection{Free energy expansion}

The free energy $F(\lambda) = -\log Z(\lambda)$ admits a cumulant expansion:
\begin{equation}\label{eq:free-energy}
    F(\lambda) = F_{\BoW} + \lambda\, \langle V \rangle_0 + \frac{\lambda^2}{2}\Big(\langle V^2 \rangle_0 - \langle V \rangle_0^2\Big) + O(\lambda^3)\,,
\end{equation}
where $\langle \cdot \rangle_0$ denotes expectation under the BoW measure~\eqref{eq:H0}.

\subsection{First-order correction}

Under the BoW measure, positions are exchangeable: $P_0(y_p {=} a,\, y_{p+1} {=} b) = P_a\, P_b$. Therefore
\begin{equation}\label{eq:first-order}
    \langle V \rangle_0 = \lambda(L-1) \sum_{a,b} P_a\, P_b\; \ee_a^{(0)} \cdot \ee_b^{(0)}\,.
\end{equation}
This is a constant independent of the token sequence: it shifts the free energy uniformly but does not affect relative probabilities or the optimal embedding. \textbf{The first-order correction to the embedding vanishes under the BoW measure.}

\subsection{Subtlety: expansion base}

The vanishing of the first-order term is a consequence of expanding around the \emph{theoretical} BoW measure, where positions are independent. If instead one fits the model to \emph{empirical data}---which has genuine adjacency correlations $P_{\text{data}}(y_p {=} a,\, y_{p+1} {=} b) \neq P_a P_b$---then the first-order term is non-trivial and captures precisely the discrepancy between bigram statistics and the bag-of-words prediction.

This raises a modeling choice: expand around the theoretical BoW measure (first non-trivial correction at $O(\lambda^2)$) or around the empirical distribution (non-trivial already at $O(\lambda)$).

\subsection{Second-order correction: the connected 4-point function}

The leading non-trivial correction under the BoW measure is
\begin{equation}\label{eq:second-order}
    \Delta F^{(2)} = \frac{\lambda^2}{2}\,\mathrm{Var}_0(V) = \frac{\lambda^2}{2}\Big(\langle V^2 \rangle_0 - \langle V \rangle_0^2\Big)\,.
\end{equation}
Expanding $V^2$:
\begin{equation}
    V^2 = \lambda^2 \sum_{p, q} (\ee_{y_p} \cdot \ee_{y_{p+1}})(\ee_{y_q} \cdot \ee_{y_{q+1}})\,.
\end{equation}
The connected part involves the 4-point correlator
\begin{equation}\label{eq:4point}
    \langle (\ee_{y_p} \cdot \ee_{y_{p+1}})(\ee_{y_q} \cdot \ee_{y_{q+1}}) \rangle_0^{\text{conn}} = \langle (\ee_{y_p} \cdot \ee_{y_{p+1}})(\ee_{y_q} \cdot \ee_{y_{q+1}}) \rangle_0 - \langle \ee_{y_p} \cdot \ee_{y_{p+1}} \rangle_0\, \langle \ee_{y_q} \cdot \ee_{y_{q+1}} \rangle_0\,,
\end{equation}
which measures the fluctuations of the nearest-neighbor correlator. Under the BoW measure, this factorizes by Wick's theorem into products of 2-point functions $\ee_a^{(0)} \cdot \ee_b^{(0)}$, giving a computable expression in terms of the unperturbed PMI matrix.

\section{Corrected embedding geometry}\label{sec:corrected}

\subsection{Perturbative expansion of the embedding}

Write the embedding as a power series in $\lambda$:
\begin{equation}\label{eq:embedding-expansion}
    \ee_i = \ee_i^{(0)} + \lambda\, \ee_i^{(1)} + \lambda^2\, \ee_i^{(2)} + \cdots
\end{equation}
The dot product expands as
\begin{equation}\label{eq:dot-expansion}
    \ee_i \cdot \ee_j = \underbrace{\ee_i^{(0)} \cdot \ee_j^{(0)}}_{=\;\Sigma_{ij}/(L P_i P_j)} + \;\lambda\underbrace{\big(\ee_i^{(0)} \cdot \ee_j^{(1)} + \ee_i^{(1)} \cdot \ee_j^{(0)}\big)}_{\delta_1(i,j)} + \;O(\lambda^2)\,.
\end{equation}

\subsection{What determines $\delta_1(i,j)$?}

The perturbation favors (or disfavors, depending on the sign of $\lambda$) sequences where adjacent tokens have large dot products. The first-order correction $\delta_1(i,j)$ should therefore be related to the \emph{adjacency PMI}:
\begin{equation}\label{eq:adj-pmi}
    \PMI_{\text{adj}}(i, j) = \log \frac{p(y_p {=} i,\, y_{p+1} {=} j)}{p(i)\, p(j)}\,,
\end{equation}
which is the PMI computed from bigram statistics rather than full-window co-occurrence. As noted in Section~\ref{sec:perturbative}, $\delta_1$ vanishes if one expands around the theoretical BoW measure but is non-trivial when fitting to empirical data.

\subsection{Self-consistency condition}

The self-referential structure implies that the correction $\ee_i^{(1)}$ depends on the unperturbed embedding $\ee^{(0)}$ through the interaction vertex. At each order in $\lambda$, the stationarity condition on the effective action yields:
\begin{equation}\label{eq:self-consistency}
    \ee_i^{(1)} \cdot \ee_j^{(0)} + \ee_i^{(0)} \cdot \ee_j^{(1)} = f\!\big(\ee^{(0)},\; \text{adjacency statistics}\big)\,,
\end{equation}
where $f$ is a function of the unperturbed embedding and the empirical bigram counts. This is a \textbf{Hartree-type self-consistency}: the correction to the propagator depends on the propagator itself.

\section{The 2PI effective action formalism}\label{sec:2PI}

\subsection{Cornwall--Jackiw--Tomboulis framework}

The natural formalism for self-consistent perturbation theory is the \textbf{CJT 2PI effective action}~\cite{cjt1974}. Define $G_{ij} \equiv \ee_i \cdot \ee_j$ as the ``propagator.'' The 2PI effective action is
\begin{equation}\label{eq:2PI}
    \Gamma[\ee, G] = \underbrace{\frac{1}{2}\mathrm{tr}\big(G_0^{-1} G\big) - \frac{1}{2}\mathrm{tr}\log G - \frac{N}{2}}_{\text{Gaussian (BoW) part}} \;+\; \underbrace{\Gamma_2[\ee, G]}_{\text{nearest-neighbor interaction}}\,,
\end{equation}
where $G_0$ is the bare propagator from the BoW theory, $G_{0,ij} = \Sigma_{ij}/(L P_i P_j)$. The interaction part collects all 2-particle-irreducible vacuum diagrams:
\begin{equation}\label{eq:Gamma2}
    \Gamma_2[\ee, G] = \lambda \sum_{p=1}^{L-1} G_{y_p, y_{p+1}} + \frac{\lambda^2}{2} \sum_{p,q} G_{y_p, y_{p+1}}\, G_{y_q, y_{q+1}}\, C_{pq}^{\text{conn}} + \cdots
\end{equation}

\subsection{The gap equation}

Physical solutions satisfy two stationarity conditions:
\begin{align}
    \frac{\delta \Gamma}{\delta \ee_i} &= 0 \qquad \text{(embedding equation of motion)}\,, \label{eq:eom-e} \\[4pt]
    \frac{\delta \Gamma}{\delta G_{ij}} &= 0 \qquad \text{(gap equation for the propagator)}\,. \label{eq:gap}
\end{align}
The gap equation~\eqref{eq:gap} determines the dressed propagator $G_{ij}$ self-consistently. At Hartree--Fock truncation (retaining only the leading diagram in $\Gamma_2$), it takes the form:
\begin{equation}\label{eq:gap-explicit}
    \boxed{\;\ee_i \cdot \ee_j = \frac{\Sigma_{ij}}{L\, P_i\, P_j} + \lambda \cdot \mathcal{A}_{ij}\big[\ee\big] + O(\lambda^2)\;\,}
\end{equation}
where $\mathcal{A}_{ij}[\ee]$ is the adjacency correction---a functional of the embedding vectors that encodes the effect of the nearest-neighbor interaction on the two-point function.

\subsection{Key properties}

The CJT formalism guarantees:
\begin{enumerate}
    \item \textbf{Thermodynamic consistency.} The dressed propagator satisfies conservation laws (Ward identities) at each truncation order.
    \item \textbf{Systematic improvement.} Higher-order corrections are organized as 2PI diagrams, each computed with the dressed propagator $G$ rather than the bare one $G_0$.
    \item \textbf{Self-consistency.} The embedding vectors $\ee_i$ appear in both the propagator ($G_{ij} = \ee_i \cdot \ee_j$) and the interaction vertex ($V \sim \lambda\, G_{y_p, y_{p+1}}$), and both are determined simultaneously by~\eqref{eq:eom-e}--\eqref{eq:gap}.
\end{enumerate}

A known subtlety: truncation of $\Gamma_2$ can introduce residual violations of global symmetries (e.g., the permutation symmetry at $\lambda = 0$ may not be recovered exactly at finite truncation order). Symmetry-improved variants of the CJT formalism~\cite{pilaftsis2013} address this.

\section{Transfer matrix perspective}\label{sec:transfer}

An alternative (non-perturbative) approach treats the model as a 1D chain and uses the transfer matrix. Define:
\begin{equation}\label{eq:transfer}
    T_{ab} = \exp\!\big(-\lambda\, \ee_a \cdot \ee_b\big)\,,
\end{equation}
where $a, b$ range over the vocabulary. The nearest-neighbor part of the partition function is
\begin{equation}\label{eq:Z-transfer}
    Z_{\text{n.n.}} = \sum_{y_1, \ldots, y_L} \prod_{p=1}^{L-1} T_{y_p, y_{p+1}} \cdot (\text{on-site weights from } H_0) = \mathrm{tr}\big(M \cdot T^{L-1}\big)\,,
\end{equation}
where $M$ encodes the BoW weights. The complication is that $H_0$ is a global function of counts, not a sum of on-site terms, so $M$ does not factorize simply.

In the perturbative regime ($\lambda \ll 1$), $T_{ab} \approx 1 - \lambda\, \ee_a \cdot \ee_b + O(\lambda^2)$, and the transfer matrix is a rank-one perturbation of the all-ones matrix (which corresponds to the BoW limit). The eigenvalue structure of $T$ is then determined by the spectrum of the Gram matrix $\ee_a \cdot \ee_b$, connecting the transfer matrix eigenvalues to the embedding geometry.

For the self-consistent problem, the transfer matrix $T$ depends on the embedding $\ee$, which depends on $T$ through the partition function---yielding a \textbf{self-consistent transfer operator}~\cite{galatolo2023}. Recent mathematical work has established existence, uniqueness, and linear response properties for such nonlinear operators.

\section{Open questions}\label{sec:open}

\begin{enumerate}
    \item \textbf{Expansion base.} Should the perturbation theory be organized around the theoretical BoW measure (where the first-order correction vanishes and the leading effect is $O(\lambda^2)$) or around the empirical distribution (where $O(\lambda)$ is non-trivial)? The former is cleaner theoretically; the latter is more natural for fitting data.

    \item \textbf{Sign of $\lambda$ and the nature of the correction.} If $\lambda > 0$ (penalizing adjacent similarity), the correction $\mathcal{A}_{ij}$ captures syntactic complementarity. If $\lambda < 0$ (rewarding adjacency), it captures semantic attraction. Can the sign be determined from data, or should both regimes be studied?

    \item \textbf{Antisymmetric structure.} The nearest-neighbor interaction distinguishes $\ee_a \cdot \ee_b$ at positions $(p, p{+}1)$ from $\ee_b \cdot \ee_a$ at $(p{+}1, p{+}2)$---i.e., it is sensitive to the \emph{direction} of adjacency. Does this induce an antisymmetric component in the corrected embedding, as in the Mallows model analysis?

    \item \textbf{Phase structure.} Is there a critical coupling $\lambda_c$ at which the permutation symmetry breaks spontaneously? In the transfer matrix language, this would correspond to a gap closing between the leading eigenvalues of $T$.

    \item \textbf{Architecture.} What neural network architecture naturally parameterizes the full model~\eqref{eq:full-H}? The BoW part is captured by Deep Sets (permutation-invariant sum pooling with quadratic readout). The nearest-neighbor part requires sensitivity to adjacent pairs---suggesting an architecture that combines sum pooling with a 1D convolutional or recurrent component, reducing to pure sum pooling at $\lambda = 0$.

    \item \textbf{Connection to the Mallows model.} The nearest-neighbor perturbation and the Mallows penalty are two different ways of breaking permutation symmetry. Are they related? In particular, does the Kendall $\tau$ penalty (which counts pairwise ordering violations) decompose into a sum of nearest-neighbor-like terms in some limit?
\end{enumerate}

\begin{thebibliography}{10}

\bibitem{cjt1974}
J.~M.~Cornwall, R.~Jackiw, and E.~Tomboulis,
``Effective potential for composite operators,''
\emph{Physical Review D}, vol.~10, pp.~2428--2445, 1974.

\bibitem{pilaftsis2013}
A.~Pilaftsis and D.~Teresi,
``Symmetry-improved CJT effective potential for the electroweak phase transition,''
\emph{Nuclear Physics B}, vol.~874, pp.~594--619, 2013.

\bibitem{galatolo2023}
S.~Galatolo,
``Mean-field coupled systems and self-consistent transfer operators: a review,''
\emph{Bollettino dell'Unione Matematica Italiana}, vol.~16, pp.~297--336, 2023.

\bibitem{metzner1995}
W.~Metzner,
``Linked-cluster expansion around the atomic limit of the Hubbard model,''
\emph{Physical Review B}, vol.~43, pp.~8549--8563, 1991.

\bibitem{arora2016}
S.~Arora, Y.~Li, Y.~Liang, T.~Ma, and A.~Risteski,
``A Latent Variable Model Approach to PMI-based Word Embeddings,''
\emph{Transactions of the Association for Computational Linguistics}, vol.~4, pp.~385--399, 2016.
arXiv:1502.03520.

\bibitem{fernandez2007}
M.~Fernandez, D.~Torres, P.~Viollier, and A.~Velasquez,
``Enertex: textual energy of associative memories for text representation,''
\emph{Lecture Notes in Computer Science}, vol.~4827, pp.~791--798, 2007.

\end{thebibliography}

\end{document}
